The Rosenbluth formula \cite{Rosenbluth:1950yq} as well as the Akhiezer-Rekalo  formalism \cite{Akhiezer:1968ek,Akhiezer:1974em}, respectively, for  unpolarized and polarized elastic $ep$ scattering hold only in the OPE formalism where two real FFs, functions of $Q^2$ only, completely describe the reaction. In the case of TPE, instead of two real functions of $Q^2$ one has to deal with three amplitudes, neglecting the electron mass. Three additional charge-odd amplitudes, that are proportional to the electron mass, are of the same order, generally of complex nature and also depending on two kinematical variables. The latter are relevant for observables such as transverse polarization.



Model-independent statements hold also in the annihilation region \cite{Gakh:2005hh,Gakh:2005wa}. Assuming time and parity invariance and crossing symmetry, the reactions $e^+e^-\leftrightarrow p\bar p$ are described by the same electromagnetic FFs. However, in the time-like region of transferred momenta FFs are of complex nature, even in the case of OPE.  Due to the large transferred momenta that are involved above the physical threshold one would therefore expect that TPE effects are enhanced. 

In the time-like region, the angular distribution of one of the produced particles allows one to extract the moduli of the FFs. In principle such measurement is simpler than a Rosenbluth separation, as it requires only one setting of the collider equipped with  a $4\pi$ detector. The $\varepsilon(\sim \cot^2\theta$)-linearity of the Rosenbluth plot translates by crossing symmetry into a $\cos^2\theta$ dependence of the annihilation cross section. Then, a forward-backward asymmetry arises in the presence of charge-odd contributions, as hard TPE. The sum (difference) of the cross section at corresponding angles would cancel (enhance) the TPE contribution, taking advantage of its charge-odd nature. 


The study of both annihilation reactions $e^+e^-\leftrightarrow p\bar p$, related by time reversal is very instructive in this respect, as radiative corrections are in general different. In $e^+e^-$  colliders, huge effort has been done to achieve a per-mille precision, based on the lepton structure function method \cite{Baier:1980kx}. This is made necessary by the use of the ``initial state radiation'' (ISR) process \cite{Czyz:2004ua}, that allows to scan a large region of transferred momenta, even in a fixed energy collision. 

Radiative correction for the reaction $\bar p p\to e^+e^- $, induced by antiproton beams  have been particularly studied \cite{Bystritskiy:2019iki} and are being implemented in a dedicated MC event generator (cf. Sec.~\ref{sec:Implementations}),  in the frame  of the time-like FF program planned at FAIR \cite{PANDA:2016fbp}. The interest of the antiproton beam  is also related to the possibility to measure the muon-antimuon final state \cite{PANDA:2020jkm}. 

The search of charge-odd contributions to the annihilation cross section has been looked for in the BaBar data, see Ref.~\cite{Tomasi-Gustafsson:2007vpn}, with the conclusion that they vanish in the limit of 2\%, which is the order of magnitude of the interference between initial and final state radiation emission. The presence of TPE in the future data from PANDA at FAIR has been simulated in Ref. \cite{PANDA:2016fbp} pointing out that a contribution of $\ge 5\% $ would be detectable. 


The QED radiative corrections to both deep-inelastic scattering and $e^+e^-$ annihilation at large virtualities
$Q^2$ and cms energies $s \gg m_e^2$ are very important for precision measurements 
since they can be 
differentially rather large in parts of the phase space. Furthermore, they strongly depend on the different 
possible measurements of the kinematic variables of the respective processes. Unlike the case in massless QCD,
the emerging logarithms $L = \ln(Q^2/m_e^2)$ or $L =\ln(s/m_e^2)$ are observable {\it phase-space} logarithms and 
cannot be resummed away as in the strictly massless case. In particular, they can be measured as has been done
in many experimental analyses in deep-inelastic scattering and $e^+e^-$ annihilation. In addition, the fermion 
mass corrections cannot be dealt with as purely massless corrections from the beginning, cf. 
e.g.~\cite{Blumlein:2020jrf}. They are, furthermore, not just given by splitting functions and massless Wilson 
coefficients, but one has to compute massive operator matrix elements (OMEs) instead of the splitting functions. 
Two 
renormalization group equations rule the corrections in the limit $Q^2 \gg m_e^2$ or $s \gg m_e^2$. \footnote{More 
generally $m_e$ can also be replaced by $m_\ell, \ell = e, \mu, \tau$} These quantities are ruled by two 
renormalization group equations, cf.~\cite{Blumlein:2000wh,Blumlein:2011mi,MUTA}, containing also the charged 
lepton masses. They read\footnote{In addition the gauge parameter $\xi$ needs to be renormalized, which we have 
suppressed here.}

\begin{align}
\label{eq:R1}
& 
\left[\left(\mu \frac{\partial}{\partial \mu} + \beta(g) \frac{\partial}{\partial g}\right) \delta_{al} 
- \gamma_m(g) m_e \frac{\partial}{\partial m_e}
+ 
\gamma_{al}(N,g) \right]\nn \\
&\times\Gamma_{li}\left(N, \frac{\mu^2}{m_e^2},g(\mu^2)\right) = 0,
\\
\label{eq:R2}
&\Biggl [\left(\mu \frac{\partial}{\partial \mu} + \beta(g) \frac{\partial}{\partial g}\right) \delta_{al} 
\delta_{kb}
- \gamma_m(g) m_e \frac{\partial}{\partial m_e}\nn \\
&- \gamma_{al}(N,g) \delta_{kb}
- \gamma_{kb}(N,g) \delta_{la} \Biggr ]\nn \\
&\times\tilde{\sigma}_{lk}\left(N, \frac{s'}{m_e^2},g(\mu^2)\right) = 0.
,\end{align}


Here $\mu$ denotes the factorization and renormalization scale, $g$ the QED coupling constant, $\beta$ the
QED $\beta$-function, $\gamma_m(g)$ the mass anomalous dimension,
$\gamma_{ij}$ the QED anomalous dimensions, $\Gamma_{li}$ are the massive OMEs and 
$\tilde{\sigma}_{lk}$ are the massless sub-system scattering coefficients. Finally, $s'$ (or $(Q^2)'$) denote the
sub-system cms energy (or virtuality). The renormalization group equations rule all corrections down to the
constant term which does not depend on the scale logarithm. Here the renormalization group equations have been written in 
Mellin space. Since the radiatively corrected (differential) scattering cross sections are observables, the 
$\mu$-dependence needs to be canceled, which is possible by a consistent expansion in the coupling constant $g$.
Usually the consideration of the radiators $\Gamma_{li}$ is therefore not sufficient, with the exception of the
leading order terms $\mathcal{O}(\alpha^k L^k)$, with $\alpha$ the fine structure constant, since there only the Born 
sub-system cross section contributes. They are leg-dependent
as the other corrections and the relevant hard scales rescale with the collinear emission variable $z$ accordingly
as will be outlined below.

The representation of the radiative corrections using Eqs.~\ref{eq:R1} and \ref{eq:R2} allows to calculate all 
but the power corrections $(m_e^2/s)^k,~~k \in \mathbb{N}, k \geq 1$. The latter terms require the full 
phase-space integrals to be carried out and are believed to yield small contributions only in the  limit of large
scales.

In Sec.~\ref{sec:2} we will consider the QED radiative corrections for 
the deep-inelastic scattering process. Section~\ref{sec:3} deals with the initial state QED radiative corrections 
to $e^+ e^-$ annihilation and Sec.~\ref{sec:4} contains the conclusions.

\subsection{Higher-order corrections to deep-inelastic scattering}\label{sec:2}


\vspace*{1mm}
\noindent
The QED corrections to deeply-inelastic scattering have been calculated in the context of different
experiments starting with those at SLAC \cite{Mo:1968cg}, CERN, with BCDMS and EMC 
\cite{Bardin:1980ii}, and with much more work for HERA, cf.~Refs.~\cite{Kwiatkowski:1990es,Arbuzov:1995id}.
These papers cover the $\mathcal{O}(\alpha)$ corrections and we will mostly discuss higher-order corrections in the 
following. Furthermore, most of the calculations refer to unpolarized-lepton unpolarized-nucleon 
scattering. Later on we will
also discuss the QED radiative corrections to polarized nucleon scattering. Early leading logarithmic  $\mathcal{O}(\alpha)$ 
approaches (LLA) are shown in Refs.~\cite{Mo:1968cg,DeRujula:1979grv,Consoli:1980nm}, see also Refs.~\cite{Kuraev:1985hb,
Kuraev:1988xn,Fadin:1989dw}. The $\mathcal{O}(\alpha)$ electroweak 
corrections, also for $\nu N$ scattering, were calculated in Refs.~\cite{Marciano:1980pb,Sarantakos:1982bp,
LlewellynSmith:1981tgo,Wheater:1982yk,Paschos:1981zx,Wirbel:1982cw,Liede:1983hw,Bardin:1986bc,Diener:2003ss,
Arbuzov:2004zr,Diener:2005me,Park:2009ft,NuTeV:2001whx} and applied to data. This has been important for the 
measurement of $\sin^2\theta_W$ from deep-inelastic data.

Complete neutral and charged current $\mathcal{O}(\alpha)$ corrections to $ep$ scattering were calculated in Refs.~\cite{Bohm:1986na,Bohm:1987cg,Spiesberger:1990fa}, calculating the bremsstrahlung contributions using MC 
integration. The neutral and charged current $\mathcal{O}(\alpha)$ corrections to deep-inelastic $ep$ scattering were
calculated analytically for the differential cross sections $\dd^2 \sigma^{\rm DIS}/\dd x/\dd Q^2$ in 
Refs.~\cite{Bardin:1987rz,Bardin:1988by,Bardin:1989vz}. The $\mathcal{O}(\alpha)$ leading log corrections to the neutral and 
charged current process in leptonic variables were computed in Ref.~\cite{Blumlein:1989gk}, correcting earlier work in Ref.~\cite{Kripfganz:1987yu}.
The LLA results agreed very well with the complete results in the neutral current case. For only leptonic 
contributions these corrections were calculated in Refs.~\cite{Beenakker:1989km,Montagna:1991ku} also. The radiative 
corrections are quite different for different choices of variables to define the differential cross sections, 
cf.~Ref.~\cite{Blumlein:2012bf}. The LLA corrections in the case of jet (or hadronic) variables were calculated 
in Ref.~\cite{Blumlein:1990wz}. In Tab.~\ref{TAB:KIN} we display the $z$ rescaling for some of the kinematic 
variables.

\begin{table*}
\centering

\caption{\label{TAB:KIN} Scaling behavior of various sets of kinematic variables for leptonic initial and final state
radiation, cf. \cite{Arbuzov:1995id,Blumlein:2012bf}. The shifted variables $\hat{Q^2}$, $\hat{y}$, the threshold $z_0$ and the Jacobian ${\cal J}$ depend on the choice of the
external kinematic variables.}
\scalebox{0.9}{
\begin{tabular}{l|c|c|c|c}

\multicolumn{5}{c}{Initial State Radiation $\hat{s} = z s$} \\
\hline
\multicolumn{1}{c}{Kinematics} &
\multicolumn{1}{|c}{$\hat{Q}^2$} &
\multicolumn{1}{|c}{$\hat{y}$} &
\multicolumn{1}{|c}{$\hat{z_0}$} &
\multicolumn{1}{|c}{${\cal J}(x,y,z)$} \\
\hline
leptonic variables
& $z Q_l^2$
& $\frac{\ds z + y_l -1}{\ds z}$
& $\frac{\ds 1- y_l }{\ds 1 - x_l y_l}$
& $\frac{\ds y_l}{\ds z + y_l -1}$
\\
mixed    variables
& $z Q_l^2$
& $\frac{\ds y_{JB}}{\ds z}$
& $y_{JB}$
& 1
\\
hadronic variables 
& $Q_h^2$
& $\frac{\ds y_h}{\ds z}$
& $y_h$
& $\frac{1}{z}$ 
\\
\hline
\multicolumn{5}{c}{Final State Radiation $\hat{s} = s$} \\
\hline
leptonic variables
& $\frac{\ds Q_l^2}{\ds z}$
& $\frac{\ds z+y_l-1}{\ds z}$
& $1 - y_l(1-x_l)$
& $\frac{\ds y_l}{\ds z(z+y_l-1)}$
\\    
mixed    variables
& $\frac{\ds Q_l^2}{\ds z}$
& $y_{JB}$
& $x_m$
& $\frac{1}{z}$
\\
\end{tabular}
}
\end{table*}

Here ${\cal J}(x,y,z)$ denotes the respective Jacobian and $x =  Q^2/(s y), y = p.q/p.l$, with 
$p$ and $l$ the nucleon and lepton 4-momentum and $q^2$ the virtuality of the process with $Q^2 = - q^2$.
$z_0$ denotes the hard bremsstrahlung threshold.
The rescaling relations for other sets of variables, such as Jaquet-Blondel variables,
the double-angle method, the $\theta y$ and  $e \Sigma$ methods, see Refs.~\cite{Arbuzov:1995id,Blumlein:2012bf}.

Let us illustrate the radiator method in lowest order LLA. One obtains

\begin{align}\frac{\dd^2 \sigma^{\rm ini(fin)}}{\dd x \dd y}& =
\frac{\alpha}{4 \pi}
L_e
\int_0^1 \dd z
 P_{ee}^{(0)}(z)  
\Biggl[
\theta(z - z_0) \nn \\
&\times{\cal J}(x,y,z)
\frac{d^2 \sigma^{0}}{\dd x \dd y}_{x=\hat{x},
y=\hat{y},s=\hat{s}}   
-  \frac{\dd^2 \sigma^{0}}{\dd x \dd y} \Biggr].
\end{align}

Here $L_e$ denotes the radiation logarithm $L_e = \ln(Q^2/m_e^2) - 1$ and we have accounted for 
a non-logarithmic correction in addition. $P_{ee}^{(0)}$ denotes the $\mathcal{O}(\alpha)$ QED splitting 
function for the $e e$ transition \cite{Fermi:1924tc}. The QED splitting functions and anomalous 
dimensions can be obtained from those of QCD \cite{Moch:2004pa,Blumlein:2021enk,
Vogt:2004mw,Ablinger:2017tan,Moch:2014sna,Blumlein:2021ryt} to three-loop order in the unpolarized and polarized cases by 
setting $T_F = 1, C_F = 1, C_A = 0$. The photon-photon splitting function needs more care.
The consistent construction of the different order radiative corrections including both the massive OMEs and the 
massless sub-system scattering cross sections is described in Ref.~\cite{Ablinger:2020qvo} 
in detail in Mellin space. Parts of the radiators depend also on generalized harmonic sums
\cite{Moch:2001zr,Ablinger:2013cf} while the others are given by harmonic sums 
\cite{Vermaseren:1998uu,Blumlein:1998if}.
However, their Mellin inversion to $z$-space can be written in terms of harmonic 
polylogarithms \cite{Remiddi:1999ew} at modified argument. Standard technologies, cf.~Ref.~\cite{Blumlein:2018cms}, 
also allow one to 
calculate the radiators in $z$-space.

We also would like to mention that the
transition matrix elements from $e^- \leftrightarrow e^+$ were given in Ref.~\cite{Blumlein:2011mi} to two-loop order. They 
are 
related to the different non-singlet anomalous dimensions $\gamma^{NS,\pm}$ \cite{Moch:2014sna,Blumlein:2021enk} 
starting at two-loop order.

Different radiative correction packages such as {\sc TERAD, HERACLES, HELIOS} and 
{\sc HECTOR} were designed  \cite{Bardin:1980ii,Kwiatkowski:1990es,Blumlein:1991ag,Akhundov:1994my,Arbuzov:1995id}. 
In addition 
to the $\mathcal{O}(\alpha L)$ corrections, which are large in certain parts of the phase space, also the $\mathcal{O}(\alpha^2 L^2)$ 
corrections were calculated for leptonic variables in Ref.~\cite{Kripfganz:1990vm} and for the whole set of kinematic 
variables in Ref.~\cite{Blumlein:1994ii}. Related radiators were obtained in Refs.~\cite{Nicrosini:1986sm,Berends:1987ab}
and later corrections in Ref.~\cite{Blumlein:2011mi}.

Likewise, complete $\mathcal{O}(\alpha)$ 
calculations have been performed for 
Jacquet-Blondel variables in Ref.~\cite{Akhundov:1993wvk} and for mixed variables in Ref.~\cite{Bardin:1995mf}. 

The 
$\mathcal{O}(\alpha^2 L)$ corrections for the latter case were computed in Ref.~\cite{Blumlein:2002fy}, extending 
the results of Refs.~\cite{Bardin:1995mf,Blumlein:1994ii}. All sub-leading corrections of this type are process 
dependent since they also depend on lower order hard scattering contributions. Therefore, resummations of 
non-leading radiators, cf.~e.g.~Ref.~\cite{Bertone:2019hks} and references therein, do not lead to a complete 
description. Furthermore, one 
has to deal with massive OMEs, cf.~e.g.~Ref.~\cite{Blumlein:2011mi}.

The QED corrections to the final state hadronic jet
are dealt with as inclusive, since they usually cannot be resolved as electromagnetic showers do also occur
in the hadronization process. Therefore, this part of the corrections is summed up according to the 
Kinoshita-Lee-Nauenberg theorem \cite{Kinoshita:1962ur,Lee:1964is}.

The radiators of the LLA corrections to higher order $\mathcal{O}(\alpha^k L^k)$ are process and radiation leg independent. Furthermore,
the non-singlet corrections are the same in the unpolarized and polarized case. The radiators for the $\mathcal{O}(\alpha^2
L^2)$ corrections were calculated in Refs.~\cite{Kripfganz:1990vm,Blumlein:1994ii}. The corresponding higher-order corrections up to $\mathcal{O}(\alpha^5 L^5)$ were computed in Refs.~\cite{Blumlein:2002bg,Jezabek:1991bx,
Skrzypek:1992vk,Przybycien:1992qe,Arbuzov:1999xso,Arbuzov:1999cq,Blumlein:2004bs,Blumlein:2007kx}, where the 
singlet corrections were calculated in Ref.~\cite{Blumlein:2004bs,Blumlein:2007kx}. One may also resum small-$x$ 
contributions of $\mathcal{O}((\alpha \ln^2(x))^k)$, which has been performed in Ref.~\cite{Blumlein:1996yz}.

In the case of leptonic variables already the $\mathcal{O}(\alpha)$ corrections do not only contain initial and final 
charged lepton radiation, but also the contributions due to the Compton peak, which was first dealt with in Ref.~\cite{Mo:1968cg} and later considered in Refs.~\cite{Bardin:1987rz,Beenakker:1989km,Blumlein:1989gk,Blumlein:1993ef}.

There are also hadronic initial state radiative corrections. In early studies 
\cite{Bohm:1986na,Bohm:1987cg,Bardin:1987rz,Bardin:1988by,Bardin:1989vz} some conceptional assumptions were still
made, which are incompatible with the massless parton model, being  generally assumed for the QCD corrections 
and
also needs to be applied to the QED corrections. The correct treatment has been given in 
Refs.~\cite{Kripfganz:1988bd,Blumlein:1990wz}. Later numerical studies were made in 
Refs.~\cite{Spiesberger:1994dm,Roth:2004ti}.

Finally, the radiative corrections to polarized-lepton polarized-nucleon scattering were also calculated at 
$\mathcal{O}(\alpha)$ in Ref.~\cite{Bardin:1996ch} and for the singlet corrections to $\mathcal{O}(\alpha^5 L^5)$ in Ref.~\cite{Blumlein:2004bs}, see also Refs.~\cite{Akushevich:1998dz,Akushevich:2011zy}.

\begin{figure}[ht]
\begin{center}
\includegraphics[width=\columnwidth,angle=0]{Higher_Order/fig/f1_hector.pdf}
\end{center}
\caption[]{
\label{epQED:FIG1}
Radiative corrections in leptonic variables in per cent for $E_e = 26.8$~GeV, $E_p =
820$~GeV. Dotted lines: $\mathcal{O}(\alpha)$, dashed lines: $\mathcal{O}(\alpha^2)$, solid lines: in addition soft
photon exponentiation, from \cite{Arbuzov:1995id}, \TCop (1995) by Elsevier Science.}
\end{figure}


Finally, we would like to give some numerical illustrations on the leptonic QED radiative corrections.
In Fig.~\ref{epQED:FIG1} we show the corrections in leptonic variables. It is shown that the LO
corrections deviate by some \% from the $\mathcal{O}(\alpha^2)$ corrections, which are therefore necessary. The soft 
exponentiation is adding much less. The reduction of the radiative corrections in deep-inelastic $ep$ scattering
is also possible by measuring the kinematic variables by other methods. Here we would like to mention, in 
particular, the double angle method, cf.~\cite{Blumlein:1994ii}, where the corrections are much smaller and do 
also behave rather flat as functions of the Bjorken variables $x$ and $y$.


\begin{figure}[ht]
\begin{center}
\includegraphics[width=\columnwidth,angle=0]{Higher_Order/fig/MIXEDa2L.pdf}
\end{center}
\caption[]{
\label{epQED:FIG2}
Radiative corrections in mixed variables in per cent for $E_e = 26.8$~GeV, $E_p =
820$~GeV. The lines left from above to below correspond to $x = 0.5, 0.1, 0.001$ in consecutive
order. Full lines $\mathcal{O}(\alpha) + \mathcal{O}(\alpha^2 L^2) +  \mathcal{O}(\alpha^2 L)$; dashed lines: $\mathcal{O}(\alpha)$;
Figure courtesy of H. Kawamura.}
\end{figure}


In Figure~\ref{epQED:FIG2} we present numerical results in the case of mixed variables 
\cite{Bardin:1995mf,Blumlein:1994ii,Blumlein:2002fy} in $\mathcal{O}(\alpha)$, corrected by the terms of $\mathcal{O}(\alpha^2 L^2)$ 
and $\mathcal{O}(\alpha^2 L)$. The dominant behaviour is given by the $\mathcal{O}(\alpha L)$ corrections. However, several \% of the 
corrections are added by the higher-order terms included. These are required for precision measurements.

\subsection{Higher-order QED corrections to \texorpdfstring{$e^+e^-$}{TEXT} annihilation}\label{sec:3}


\vspace*{1mm}
\noindent
Unlike the case in deep-inelastic $ep$ scattering, the complete two-loop QED corrections at large cms energy 
have been calculated for $e^+ e^-$ annihilation \cite{Blumlein:2020jrf}, 
which are of central importance for precision measurements of the $Z^0$ peak, $Z^0 H^0$ production and $t\bar{t}$ 
production at threshold, cf.~\cite{Blumlein:2019pqb} and references therein.\footnote{For a recent survey see 
Ref.~\cite{Blumlein:2022mrp}. An excellent survey on the different calculations at one-loop order is given in 
Ref.~\cite{BP}.}
In earlier attempts to 
calculate these corrections or parts of it in Refs.~\cite{Berends:1987ab,Kniehl:1988id} the massless limit $m_e 
\rightarrow 0$ has been taken too early, which led to incorrect results. Using for the inclusive cross 
sections the factorization property \cite{Buza:1995ie}, one may also calculate it referring to the massless 
sub-processes \cite{Hamberg:1990np,Harlander:2002wh,Blumlein:2019srk} and the corresponding massive OMEs calculated in Ref.~\cite{Blumlein:2011mi}. Since the results of Ref.~\cite{Blumlein:2011mi} did not agree 
with those from Refs.~\cite{Berends:1987ab,Kniehl:1988id}, there was also the possibility that factorization does not 
hold in the massive case. However, it could be proven by the full massive phase-space calculation in Ref.~\cite{Blumlein:2020jrf} that this is not the case and the reasons for the differences could be given in Refs.~\cite{Blumlein:2020jrf,Blumlein:2019srk,Blumlein:2019pqb}. With the factorization relation at hand one can now 
calculate the three most leading logarithms in any order of $\alpha$. This has been done to $\mathcal{O}(\alpha^6 L^5)$ in 
Ref.~\cite{Ablinger:2020qvo}. 

\begin{table*}
    \centering
    \caption{ Shifts in the $Z$-mass and the width due to the 
different contributions to the ISR QED
    radiative corrections for a fixed width of $\Gamma_Z =  
2.4952~\GeV$  and $s$-dependent width using
    $M_Z = 91.1876~\GeV$
    \cite{ParticleDataGroup:2018ovx} and $s_0 = 4 m_\tau^2$, 
cf.~\cite{ALEPH:2005ab}; from 
\cite{Ablinger:2020qvo}, \TCop (2020) by Elsevier Science.}
\scalebox{0.9}{
    \begin{tabular}{l|r|r|r|r}
    \multicolumn{1}{c}{} &
    \multicolumn{2}{|c|}{Fixed width} &
    \multicolumn{2}{c}{$s$-dep. width} \\
    \hline
    \multicolumn{1}{c|}{Order of correction} &
    \multicolumn{1}{c|}{Peak} &
    \multicolumn{1}{c|}{Width}    &
    \multicolumn{1}{c|}{Peak} &
    \multicolumn{1}{c}{Width} \\
    \multicolumn{1}{c|}{} &
    \multicolumn{1}{c|}{(MeV)} &
    \multicolumn{1}{c|}{(MeV)}    &
    \multicolumn{1}{c|}{(MeV)} &
    \multicolumn{1}{c}{(MeV)} \\
\hline
    %lowest order                            &       &     &           &          \\
    $\mathcal{O}(\alpha)$                    &   185.638   &   
539.408  &   181.098  &   524.978 \\
    $\mathcal{O}(\alpha^2 L^2)$                       & - 96.894   & 
-177.147  & - 95.342  & -176.235 \\
    $\mathcal{O}(\alpha^2 L)$                         &     6.982   &    
22.695  &     6.841  &    21.896 \\
    $\mathcal{O}(\alpha^2 )$                          &     0.176   & -  
2.218  &     0.174  & -  2.001 \\
    $\mathcal{O}(\alpha^3 L^3)$                       &     23.265  &    
38.560  &    22.968  &    38.081 \\
    $\mathcal{O}(\alpha^3 L^2)$                       & -   1.507  & -  
1.888  & -  1.491  & -  1.881 \\
    $\mathcal{O}(\alpha^3 L)$                        & -   0.152  &     
0.105  & -  0.151  & -  0.084 \\
    $\mathcal{O}(\alpha^4 L^4)$                       & -   1.857  &     
0.206  & -  1.858  &     0.146 \\
    $\mathcal{O}(\alpha^4 L^3)$                       &      0.131  & -  
0.071  &     0.132  & -  0.065 \\
    $\mathcal{O}(\alpha^4 L^2)$                       &      0.048  & -  
0.001  &     0.048  &     0.001 \\
    $\mathcal{O}(\alpha^5 L^5)$                       &      0.142  & -  
0.218  &     0.144  & -  0.212 \\
    $\mathcal{O}(\alpha^5 L^4)$                       & -   0.000  &     
0.020  & -  0.001  &     0.020 \\
    $\mathcal{O}(\alpha^5 L^3)$                       & -   0.008  &     
0.009  & -  0.008  &     0.008 \\
    $\mathcal{O}(\alpha^6 L^6)$                       & -   0.007  &     
0.027  & -  0.007  &     0.027 \\
    $\mathcal{O}(\alpha^6 L^5)$                       & -   0.001  &     
0.000  & -  0.001  &     0.000 \\
    \end{tabular}}
    \label{TAB1}
    \end{table*}

Besides the two-loop massive OME $A_{ee}^{(2)}$ \cite{Blumlein:2011mi} one also needs the two-loop massive OME 
$A_{\gamma e}^{(2)}$ \cite{Ablinger:2020qvo}

    \begin{align}
A_{\gamma e} &= a^{\sf \overline{MS}}
\Biggl[-\frac{1}{2} P_{\gamma e}^{(0)} L_1 
+ \Gamma_{\gamma e}^{(0)}\Biggr] \nn\\
& + {a^{\sf \overline{MS}}}^2
\Biggl[
\frac{P_{\gamma e}^{(0)}}{8} (  
P_{e e}^{(0)} +  
P_{\gamma \gamma}^{(0)} + 2 \beta_0) L_1^2
\nonumber\\ &
+ \frac{1}{2} ( P_{\gamma e}^{(1)} 
+ \Gamma_{e e}^{(0)}  P_{\gamma e}^{(0)} 
+ \Gamma_{\gamma e}^{(0)}  P_{\gamma \gamma}^{(0)}  
\nn \\
&+ 2 \beta_0 \Gamma_{\gamma e}^{(0)} ) L_1   
+ \hat{\Gamma}_{\gamma e}^{(1)}  
+ \bar{\Gamma}_{e e}^{(0)} P_{\gamma e}^{(0)}  
\nonumber\\ &
+ \bar{\Gamma}_{\gamma e}^{(0)} P_{\gamma \gamma}^{(0)}  
+ 2 \beta_0 \bar{P}_{\gamma e}^{(0)} 
\Biggr],
\end{align}


\noindent
with $L_1 = \ln(m^2/\mu^2)$, $\hat{\Gamma}_{ij} =
\Gamma_{ij}(N_F+1) - \Gamma_{ij}(N_F)$ and $\bar{f}$ are the 
$\mathcal{O}(\epsilon)$ parts of the functions $f$, where $\epsilon = 2 - D/2$ denotes 
the dimensional parameter. 

In Tab.~\ref{TAB1} we show the improvement of the peak position and the size of the width of the 
$Z^0$ boson by including higher and higher order QED initial state corrections, until we reach the 
level of $\pm 10$~keV as the theoretical error. We consider both the cases of an $s$-independent 
as well of the $s$-dependent width $\Gamma_Z$.

The QED initial state corrections are also very important for the measurement of inclusive 
$t\bar{t}$ production in the threshold region, which will be used for a precision measurement of 
the top quark mass in the future. Evidently higher than the $\mathcal{O}(\alpha^2)$ corrections are required
to obtain a sufficient description, cf.~Fig.~\ref{fig:TT1}.
%%%%%%%%%%%%%%%%%%%%%%%%%%%%%%%%%%%%%%%%%%%%%%%%%%%%%%%%%%%%%%%%%%%%%%%%
\begin{figure}[ht]
  \centering
  \hskip-0.8cm
  \includegraphics[width=\columnwidth]{Higher_Order/fig/TTbarTHR.pdf}
  \caption{ The QED  ISR corrections to $e^+e^- \rightarrow 
t\overline{t}$ ($s$-channel photon exchange) in the threshold
  region for a PS-mass of $m_t = 172~\GeV$.
Dotted line $\mathcal{O}(\alpha^{0})$;
Dashed line $\mathcal{O}(\alpha)$;
Dash-dotted line $\mathcal{O}(\alpha^{2})$;
Full line $\mathcal{O}(\alpha^{2})$ + soft resummation; from 
Ref.~\cite{Blumlein:2019pqb}, \TCop (2019) by Elsevier Science.
\label{fig:TT1}}
\end{figure}   
%%%%%%%%%%%%%%%%%%%%%%%%%%%%%%%%%%%%%%%%%%%%%%%%%%%%%%%%%%%%%%%%%%%%%%%%

For a future precision measurement of the fine structure constant $\alpha$ at high 
energies \cite{Janot:2015gjr} one needs also higher-order corrections to the forward-backward asymmetry. The 
corresponding LO corrections to $\mathcal{O}(\alpha^6 L^6)$ have been calculated in Ref.~\cite{Blumlein:2021jdl}
extending earlier results at $\mathcal{O}(\alpha^2 L^2)$ \cite{Beenakker:1989km}. The 
$\mathcal{O}(\alpha)$ corrections were computed in Ref.~\cite{Bohm:1989pb}.

The LLA initial state radiative corrected  forward-backward asymmetry is 
given by

\begin{eqnarray}
A_{\rm FB}(s) &=& \frac{1}{\sigma_T(s)} \int_{z_0}^1 \dd z 
\frac{4z}{(1+z)^2} 
\nonumber\\ && \times
\tilde{H}_{e}^{\rm LLA}(z) \sigma_{\rm FB}^{(0)}(zs),
\end{eqnarray}
with $\sigma_{\rm FB}^{(0)} = \sigma_{\rm F}^{(0)} - \sigma_{\rm B}^{(0)}$ and 
$\sigma_T(s)$ the radiatively corrected total cross section. The radiator is given 
by

\begin{eqnarray}
\tilde{H}_e^{\rm LLA}(z) =
\left[ H_{e}^{\rm LLA}(z)
+ H_{\rm FB}^{\rm LLA}(z)  \right],
\end{eqnarray}
where the parameter $z_0$ plays the role of an energy cut and $z = 
s'/s$. Here $H_{e}^{\rm LLA}$ is angular independent, and
denotes the leading logarithmic contributions of the usual radiators, while 
$H_{\rm FB}^{\rm LLA}$, which is 
angular dependent and obtained by the following integral


\begin{align}        H_{\rm FB}^{\rm LLA}(z) &= \int\limits_{0}^{1} {\rm d} x_1 
\int\limits_{0}^{1} {\rm d} x_2 \Biggl [ \left(
\frac{(1+z)^2}{(x_1+x_2)^2} -
        1 \right) \nn\\
        &\times\Gamma_{ee}^{\rm LLA}(x_2) \Gamma_{ee}^{\rm LLA}(x_1) 
\delta(x_1 x_2 - z) \Biggr ]. \label{eq:conv}
\end{align}

Here $\Gamma_{ee}^{\rm LLA}(x_i)$ denote the {massive OMEs} and (\ref{eq:conv}) is expanded in
$(\alpha L)$ to the respective order needed. One may use the Mellin-transform

\begin{align} 
        &\mathcal{M}[H_{\rm FB}^{\rm LLA}(z)](n) = \int\limits_{0}^{1} {\rm 
d}z z^n H_{\rm FB}^{\rm LLA}(z) \nonumber\\
&=
\int\limits_{0}^{1} {\rm d}x_1
        \int\limits_{0}^{1} {\rm d}x_2 x_1^n x_2^n \Biggl [ \left( 
\frac{(1+x_1x_2)^2}{(x_1+x_2)^2} - 1 \right) \nn \\   
&\times\Gamma_{ee}^{\rm LLA}(x_2)
        \Gamma_{ee}^{\rm LLA}(x_1) \Biggr ].
\end{align}

In $z$-space the radiator depends on cyclotomic harmonic polylogarithms 
\cite{Ablinger:2011te}.
%\renewcommand{\arraystretch}{1.2}
\begin{table}[ht]
   \centering
       \caption{ $A_{\rm FB}$ evaluated at $s_-=(87.9\, {\rm GeV})^2$, $M_Z^2$
                        and $s_+=(94.3\,{\rm GeV})^2$ for the cut $z>4m_\tau^2/s$ from
        Ref.~\cite{Blumlein:2021jdl}, \TCop (2021) by Elsevier Science.}
        \label{tab:AFB1}
    \begin{tabular}{r|r|r|r}
    \multicolumn{1}{c|}{ Order} &
    \multicolumn{1}{c|}{$A_{\rm FB}(s_-)$  } &
    \multicolumn{1}{c|}{$A_{\rm FB}(M_Z^2)$} &
    \multicolumn{1}{c}{$A_{\rm FB}(s_+)$} \\
    \hline
                $\mathcal{O}(\alpha^0)$         & $-0.3564803$          
& $ 0.0225199$  & $0.2052045$
                \\
                $\mathcal{O}(\alpha L^1)$       & $-0.2945381$          
& $-0.0094232$  & $0.1579347$
                \\
                $\mathcal{O}(\alpha L^0)$       & $-0.2994478$          
& $-0.0079610$  & $0.1611962$
                \\
                $\mathcal{O}(\alpha^2 L^2)$     & $-0.3088363$          
& $ 0.0014514$  & $0.1616887$
                \\
                $\mathcal{O}(\alpha^3 L^3)$     & $-0.3080578$          
& $ 0.0000198$  & $0.1627252$
                \\
                $\mathcal{O}(\alpha^4 L^4)$     & $-0.3080976$          
& $ 0.0001587$  & $0.1625835$
                \\
                $\mathcal{O}(\alpha^5 L^5)$     & $-0.3080960$          
& $ 0.0001495$  & $0.1625911$
                \\
                $\mathcal{O}(\alpha^6 L^6)$     & $-0.3080960$          
& $ 0.0001499$  & $0.1625911$
\\
        \end{tabular}
\end{table}

\renewcommand{\arraystretch}{1.0}

\noindent
In Tab.~\ref{tab:AFB1} we illustrate the improvement of the forward-backward 
asymmetry by including higher order leading-log initial state radiative 
corrections beyond the $\mathcal{O}(\alpha^2 L^2)$ corrections \cite{Beenakker:1989km}. The 
$\mathcal{O}(\alpha^6 L^6)$ corrections stabilize the off-resonance forward-backward asymmetries to 
six digits. Of course, also non-leading log corrections are needed to be 
calculated in the future.

\subsection{Conclusions}\label{sec:4}

\vspace*{1mm}
\noindent
We have summarized the status of the higher-order QED corrections to deep-inelastic $ep$ scattering and 
for the QED initial state corrections to $e^+e^-$ annihilation. In both cases the emphasis is on the non-power 
corrections, i.e.~on the logarithmic contributions down to the constant terms of $\mathcal{O}(\alpha^k L^0)$. These 
contributions can be computed using renormalization group techniques, which involve massive OMEs on one side 
and massless sub-system scattering processes on the other side. It is important in using this approach to cancel
the dependence on the factorization scale by matching the different contributions. This is best done in 
Mellin-$N$ space in a fully analytic way. The radiators are then given by Mellin convolutions of the respective 
pieces defining the OMEs and the sub-system scattering cross sections. The result can be Mellin-inverted
analytically providing the correct radiators. The requested interplay between both types of contributions does not
allow to study QED evolution only. 

In the case of future experimental studies of $e^+e^-$ annihilation it is planned to measure the width and the 
mass of the $Z^0$-boson down to the $20$~keV range. This requires initial state corrections of up to 6th order
QED corrections by keeping three logarithmic orders.

